\makechapterstyle{gradualesimplex}{%
  % Define o espaçamento anterior ao número do capítulo
  % (1) baselineskip + AlturaCabecalhoExtra = OK
  % (2) baselineskip + AlturaCabecalho = INSUFICIENTE
  % (3) AlturaCabecalhoExtra + AlturaCabecalho = similar a (1)
  % (4) baselineskip + AlturaCabecalhoExtra + AlturaCabecalho = EXAGERADO
  \setlength{\beforechapskip}{-\baselineskip}
  \addtolength{\beforechapskip}{-\AlturaCabecalho}
  % Define o espaçamento entre o número e o título do capítulo
  \setlength{\midchapskip}{-\baselineskip}
  % Define o espaçamento posterior ao título do capítulo
  \setlength{\afterchapskip}{\baselineskip}
  %\renewcommand{\chapnamefont}{}
  %\renewcommand{\chapnumfont}{}
  \renewcommand{\booktitlefont}{\scshape\HUGE\centering}
  \renewcommand{\printbookname}{}
  \renewcommand{\printbooknum}{}
  \renewcommand{\afterbookskip}{\relax}
  \renewcommand{\parttitlefont}{\scshape\LARGE\centering}
  \renewcommand{\printpartname}{}
  \renewcommand{\printpartnum}{}
  \renewcommand{\afterpartskip}{\vspace{2.5\baselineskip}}
  \renewcommand{\chaptitlefont}{\bfseries\Large\centering}
  \renewcommand{\printchaptername}{}
  \renewcommand{\printchapternum}{}
  % Define o estilo das seções
  \setsecheadstyle{\bfseries\large\centering}
  % Omite a numeração das seções
  %\renewcommand{\thesection}{}
  % Define o estilo das subseções
  \setsubsecheadstyle{\bfseries\centering}
  \setsubsubsecheadstyle{\bfseries\small\centering}
  % Omite a numeração das subseções
  %\renewcommand{\thesubsection}{}
}

\NewDocumentCommand{\ubook}{m}{%
  \openright
  \book*{#1}
  \phantomsection
  \addcontentsline{toc}{book}{#1}
}

\NewDocumentCommand{\upart}{m}{%
  \nopartblankpage
  \begingroup
  \let\cleardoublepage\relax
  \part*{#1}
  \endgroup
  \phantomsection
  \addcontentsline{toc}{part}{#1}
  \openany
}

\NewDocumentCommand{\uchapter}{s O{#3} m O{#3}}{%
  \begingroup
  % Evita quebra de página antes do início do capítulo
  \IfBooleanT{#1}{\let\clearpage\relax}
  \chapter*{#3}
  \endgroup
  \phantomsection
  \addcontentsline{toc}{chapter}{#2}
  \chaptermark{#4}
}

\NewDocumentCommand{\usection}{O{#2} m O{#2}}{%
  \section*{#2}
  \phantomsection
  \addcontentsline{toc}{section}{#1}
  \sectionmark{#3}
}

\NewDocumentCommand{\usubsection}{m}{%
  \subsection*{#1}
  \phantomsection
  \addcontentsline{toc}{subsection}{#1}
}

\NewDocumentCommand{\uchaptercomplex}{s O{#3} m O{#3} m O{}}{%
\IfBooleanTF{#1}{\uchapter*}{\uchapter}[#2]{%
{\normalfont\normalsize\color{gregoriocolor}#5} \\[-6pt]
#3 \\[-6pt]
{\normalfont\normalsize\color{gregoriocolor}#6}%
}[#4]
}

\NewDocumentCommand{\Solemnitas}{s O{#3} m O{#3} m}{%
  \IfBooleanTF{#1}{\uchaptercomplex*}{\uchaptercomplex}[#2]{#3}[#4]{#5}[Solenidade]%
}

\NewDocumentCommand{\Festum}{s O{#3} m O{#3} m}{%
\IfBooleanTF{#1}{\uchaptercomplex*}{\uchaptercomplex}[#2]{#3}[#4]{#5}[Festa]%
}